%% College Assignment Collection starter -- adapted from bimalka98's
%% "Assignment Template - 2 Pages Compact" layout.
%%
%% The upstream template ships figures as raster images
%% (\includegraphics{figures/...}); this starter draws its one figure
%% with TikZ instead, so the download is fully self-contained with no
%% missing-file risk.
%%
%% Source: https://github.com/bimalka98/LaTeX-Templates
%% (Assignment_Template_2_Pages_Compact/)
%% Licence: MIT, (c) 2021 Bimalka Piyaruwan.
\documentclass[a4paper,11pt,twocolumn]{article}

%% packages

\usepackage{amsmath} % needed for command eqref
\usepackage{amssymb} % needed for math fonts
\usepackage[colorlinks=true,breaklinks]{hyperref} % creates hyperlinks; colorlinks=true removes the boxes, breaklinks breaks long links
\usepackage[hyphenbreaks]{breakurl} % breaks long URLs in bibliography references, including across a hyphenated line
\usepackage{xcolor}
\definecolor{c1}{rgb}{0,0,1} % blue
\definecolor{c2}{rgb}{0,0.3,0.9} % light blue
\definecolor{c3}{rgb}{0.3,0,0.9} % red-blue
\hypersetup{
    linkcolor={c1}, % internal links
    citecolor={c2}, % citations
    urlcolor={c3} % external links/urls
}
\usepackage[square,authoryear]{natbib} % needed for cite and abbrvnat bibliography style
\usepackage[nottoc]{tocbibind} % needed for displaying bibliography and other in the table of contents
\usepackage{graphicx} % needed for \includegraphics
\usepackage{tikz} % needed for the vector diagram in this starter
\usepackage{longtable} % needed for long tables over pages
\usepackage{bigstrut} % needed for the command \bigstrut
\usepackage{enumerate} % needed for some options in enumerate
\usepackage{makeidx} % needed for creating an index
\makeindex
\usepackage{gensymb}
\usepackage{url}

%% page settings

\usepackage[top=5mm, bottom=15mm,left=8mm,right=8mm]{geometry} % needed for page border settings
\parindent=0mm % for space of first line of new text block
\sloppy % for writing with hyphenless justification (tries to)
\hyphenation{} % use hyphenation of tolerance parametershttp://www.jr-x.de/publikationen/latex/tipps/zeilenumbruch.html
\hyphenpenalty=10000
\exhyphenpenalty=10000
\usepackage{fancyhdr} % needed for head and foot options
%% my macros

%% Text fomats
\newcommand{\tbi}[1]{\textbf{\textit{#1}}}

%% Math fonts
\newcommand{\bbA}{\mathbb{A}}
\newcommand{\bbB}{\mathbb{B}}
\newcommand{\bbC}{\mathbb{C}}
\newcommand{\bbD}{\mathbb{D}}
\newcommand{\bbE}{\mathbb{E}}
\newcommand{\bbF}{\mathbb{F}}
\newcommand{\bbG}{\mathbb{G}}
\newcommand{\bbH}{\mathbb{H}}
\newcommand{\bbI}{\mathbb{I}}
\newcommand{\bbJ}{\mathbb{J}}
\newcommand{\bbK}{\mathbb{K}}
\newcommand{\bbL}{\mathbb{L}}
\newcommand{\bbM}{\mathbb{M}}
\newcommand{\bbN}{\mathbb{N}}
\newcommand{\bbO}{\mathbb{O}}
\newcommand{\bbP}{\mathbb{P}}
\newcommand{\bbQ}{\mathbb{Q}}
\newcommand{\bbR}{\mathbb{R}}
\newcommand{\bbS}{\mathbb{S}}
\newcommand{\bbT}{\mathbb{T}}
\newcommand{\bbU}{\mathbb{U}}
\newcommand{\bbV}{\mathbb{V}}
\newcommand{\bbW}{\mathbb{W}}
\newcommand{\bbX}{\mathbb{X}}
\newcommand{\bbY}{\mathbb{Y}}
\newcommand{\bbZ}{\mathbb{Z}}

\title{\Large\textbf{Design and Analysis of a Low-Pass RC Filter for
  Sensor Signal Conditioning}}
\author{Ada Lovelace\hspace{1cm}20270001}
\date{\today}

\begin{document}
\maketitle

A first-order RC low-pass filter is a common building block for
removing high-frequency noise from an analog sensor signal before it
reaches an analog-to-digital converter. This report derives its
transfer function, selects component values for a target cutoff
frequency, and verifies the design against a measured noise spectrum.

\begin{center}
\begin{figure}[!h]
\centering
\begin{tikzpicture}[scale=1, every node/.style={font=\small}]
  % input node
  \draw (0,0) node[left] {$V_{\text{in}}$} -- (0.6,0);
  % resistor (zig-zag)
  \draw (0.6,0) -- (0.9,0.15) -- (1.1,-0.15) -- (1.3,0.15)
        -- (1.5,-0.15) -- (1.7,0.15) -- (2.0,0) -- (2.6,0);
  \node[above] at (1.3,0.15) {$R$};
  % node to output
  \draw (2.6,0) -- (3.4,0) node[right] {$V_{\text{out}}$};
  \filldraw (2.6,0) circle (1pt);
  % capacitor to ground
  \draw (2.6,0) -- (2.6,-0.4);
  \draw (2.4,-0.4) -- (2.8,-0.4);
  \draw (2.4,-0.55) -- (2.8,-0.55);
  \node[right] at (2.8,-0.475) {$C$};
  \draw (2.6,-0.55) -- (2.6,-0.9);
  \draw (2.4,-0.9) -- (2.8,-0.9);
  \draw (2.45,-1.0) -- (2.75,-1.0);
  \draw (2.5,-1.1) -- (2.7,-1.1);
\end{tikzpicture}
\caption{First-order RC low-pass filter: resistor $R$ in series,
  capacitor $C$ to ground at the output node.}
\end{figure}
\end{center}

\section{Transfer Function}

Treating the capacitor's impedance as $Z_C = 1/(j\omega C)$, the
voltage divider formed by $R$ and $C$ gives
\begin{equation}
  H(j\omega) = \frac{V_{\text{out}}}{V_{\text{in}}}
  = \frac{1/(j\omega C)}{R + 1/(j\omega C)}
  = \frac{1}{1 + j\omega RC}.
  \label{eq:transfer}
\end{equation}
The magnitude response is $|H(j\omega)| = 1/\sqrt{1 + (\omega RC)^2}$,
which is $1/\sqrt{2}$ (the $-3\,\text{dB}$ point) when
$\omega RC = 1$, giving the cutoff frequency
\begin{equation}
  f_c = \frac{1}{2\pi RC}.
  \label{eq:cutoff}
\end{equation}

\section{Component Selection}

The sensor noise floor was measured to be dominated by components above
$500\,\text{Hz}$, so a cutoff of $f_c = 200\,\text{Hz}$ was chosen to
give attenuation margin without cutting into the signal band (DC to
$50\,\text{Hz}$). Fixing $C = 100\,\text{nF}$ (a common, readily
available value) and solving \eqref{eq:cutoff} for $R$:
\begin{equation}
  R = \frac{1}{2\pi f_c C}
  = \frac{1}{2\pi \cdot 200 \cdot 100 \times 10^{-9}}
  \approx 7.96\,\text{k}\Omega,
\end{equation}
rounded to the nearest standard value, $R = 8.2\,\text{k}\Omega$
(E12 series), giving an actual cutoff of $f_c \approx 194\,\text{Hz}$.

\section{Verification}

Table~\ref{tab:atten} lists the predicted attenuation at several test
frequencies from \eqref{eq:transfer}, using the standard resistor value.

\begin{table}[!h]
\centering
\caption{Predicted attenuation at test frequencies.}
\label{tab:atten}
\begin{tabular}{cc}
\hline
Frequency (Hz) & Attenuation (dB) \\
\hline
50   & $-1.1$ \\
194  & $-3.0$ \\
500  & $-8.4$ \\
1000 & $-14.2$ \\
\hline
\end{tabular}
\end{table}

The signal band (up to $50\,\text{Hz}$) loses less than $1.2\,\text{dB}$,
while noise above $500\,\text{Hz}$ is attenuated by more than
$8\,\text{dB}$, consistent with the design goal.

\section{Conclusion}

A single-stage RC low-pass filter with $R = 8.2\,\text{k}\Omega$ and
$C = 100\,\text{nF}$ meets the target cutoff of approximately
$200\,\text{Hz}$ while preserving the signal band, using only two
low-cost, readily available components \citep{horowitz2015art}.

\bibliographystyle{abbrvnat}
\bibliography{reference}

\end{document}
